\section{Mat\-Mgnt3d Class Reference}
\label{classMatMgnt3d}\index{MatMgnt3d@{MatMgnt3d}}
{\tt \#include $<$matmgnt3d.hpp$>$}



\subsection{Detailed Description}




Definition at line 33 of file matmgnt3d.hpp.\subsection*{Public Types}
\begin{CompactItemize}
\item 
enum \{ {\bf UE} = 0, 
{\bf UC} = 1, 
{\bf DE} = 2, 
{\bf DC} = 3
 \}
\end{CompactItemize}
\subsection*{Public Member Functions}
\begin{CompactItemize}
\item 
{\bf Mat\-Mgnt3d} ()
\item 
{\bf $\sim$Mat\-Mgnt3d} ()
\item 
{\bf Kernel3d} \& {\bf knl} ()
\item 
int \& {\bf np} ()
\item 
double {\bf alt} ()
\item 
int {\bf src\-DOF} ()
\item 
int {\bf trg\-DOF} ()
\item 
int {\bf dim} ()
\item 
int {\bf setup} ()
\item 
int {\bf report} ()
\item 
int {\bf pln\-Dat\-Sze} (int tp)
\item 
int {\bf eff\-Dat\-Sze} (int tp)
\item 
int {\bf Upw\-Chk2Upw\-Equ\_\-dgemv} (int level, const Dbl\-Num\-Vec \&, Dbl\-Num\-Vec \&)
\item 
int {\bf Upw\-Equ2Upw\-Chk\_\-dgemv} (int level, Index3 ii, const Dbl\-Num\-Vec \&, Dbl\-Num\-Vec \&)
\item 
int {\bf Dwn\-Chk2Dwn\-Equ\_\-dgemv} (int level, const Dbl\-Num\-Vec \&, Dbl\-Num\-Vec \&)
\item 
int {\bf Dwn\-Equ2Dwn\-Chk\_\-dgemv} (int level, Index3 ii, const Dbl\-Num\-Vec \&, Dbl\-Num\-Vec \&)
\item 
int {\bf Upw\-Equ2Dwn\-Chk\_\-dgemv} (int level, Index3 ii, const Dbl\-Num\-Vec \&eff\-Den, Dbl\-Num\-Vec \&eff\-Val)
\item 
int {\bf pln\-Den2Eff\-Den} (int level, const Dbl\-Num\-Vec \&, Dbl\-Num\-Vec \&)
\begin{CompactList}\small\item\em plain-$>$regular-$>$effective \item\end{CompactList}\item 
int {\bf sam\-Den2Reg\-Den} (const Dbl\-Num\-Vec \&, Dbl\-Num\-Vec \&)
\item 
int {\bf eff\-Val2Pln\-Val} (int level, const Dbl\-Num\-Vec \&, Dbl\-Num\-Vec \&)
\begin{CompactList}\small\item\em effective-$>$regular-$>$plain \item\end{CompactList}\item 
int {\bf reg\-Val2Sam\-Val} (const Dbl\-Num\-Vec \&, Dbl\-Num\-Vec \&)
\item 
const Dbl\-Num\-Mat \& {\bf sam\-Pos} (int tp)
\item 
const Dbl\-Num\-Mat \& {\bf reg\-Pos} ()
\item 
int {\bf local\-Pos} (int, Point3, double, Dbl\-Num\-Mat \&)
\item 
int {\bf sam\-Pos\-Cal} (int n, double R, Dbl\-Num\-Mat \&ret)
\item 
int {\bf reg\-Pos\-Cal} (int n, double R, Dbl\-Num\-Mat \&ret)
\item 
int {\bf cptwvv} (int, double, fftw\_\-complex $\ast$, int, fftw\_\-complex $\ast$, int, fftw\_\-complex $\ast$, int)
\end{CompactItemize}
\subsection*{Static Public Member Functions}
\begin{CompactItemize}
\item 
static {\bf Mat\-Mgnt3d} $\ast$ {\bf getmmptr} ({\bf Kernel3d}, int)
\end{CompactItemize}
\subsection*{Public Attributes}
\begin{CompactItemize}
\item 
{\bf Kernel3d} {\bf \_\-knl}
\item 
int {\bf \_\-np}
\item 
bool {\bf \_\-hom}
\item 
vector$<$ double $>$ {\bf \_\-deg\-Vec}
\item 
map$<$ int, Dbl\-Num\-Mat $>$ {\bf \_\-upw\-Chk2Upw\-Equ}
\item 
map$<$ int, Num\-Tns$<$ Dbl\-Num\-Mat $>$ $>$ {\bf \_\-upw\-Equ2Upw\-Chk}
\item 
map$<$ int, Dbl\-Num\-Mat $>$ {\bf \_\-dwn\-Chk2Dwn\-Equ}
\item 
map$<$ int, Num\-Tns$<$ Dbl\-Num\-Mat $>$ $>$ {\bf \_\-dwn\-Equ2Dwn\-Chk}
\item 
map$<$ int, Off\-Tns$<$ Dbl\-Num\-Mat $>$ $>$ {\bf \_\-upw\-Equ2Dwn\-Chk}
\item 
Dbl\-Num\-Mat {\bf \_\-sam\-Pos} [4]
\item 
Dbl\-Num\-Mat {\bf \_\-reg\-Pos}
\end{CompactItemize}
\subsection*{Static Public Attributes}
\begin{CompactItemize}
\item 
static double {\bf \_\-wsbuf} [$\,$]
\item 
static vector$<$ {\bf Mat\-Mgnt3d} $>$ {\bf \_\-mmvec}
\end{CompactItemize}


\subsection{Member Enumeration Documentation}
\subsubsection{\setlength{\rightskip}{0pt plus 5cm}anonymous enum}\label{classMatMgnt3d_dca29a1140aadadfd92b34a02fa516ef_dca29a1140aadadfd92b34a02fa516ef}


UE = Upper Equivalent UC = Upper Check DE = Downward Equivalent DC = Downward Check \begin{Desc}
\item[Enumerator: ]\par
\begin{description}
\index{UE@{UE}!MatMgnt3d@{MatMgnt3d}}\index{MatMgnt3d@{MatMgnt3d}!UE@{UE}}\item[{\em 
UE\label{classMatMgnt3d_dca29a1140aadadfd92b34a02fa516ef_dca29a1140aadadfd92b34a02fa516ef0612e93d65abcdaeed533a0876474251}
}]\index{UC@{UC}!MatMgnt3d@{MatMgnt3d}}\index{MatMgnt3d@{MatMgnt3d}!UC@{UC}}\item[{\em 
UC\label{classMatMgnt3d_dca29a1140aadadfd92b34a02fa516ef_dca29a1140aadadfd92b34a02fa516ef6a2fb30a2dacdda4d74122697262a8b7}
}]\index{DE@{DE}!MatMgnt3d@{MatMgnt3d}}\index{MatMgnt3d@{MatMgnt3d}!DE@{DE}}\item[{\em 
DE\label{classMatMgnt3d_dca29a1140aadadfd92b34a02fa516ef_dca29a1140aadadfd92b34a02fa516ef3a52f3c22ed6fcde5bf696a6c02c9e73}
}]\index{DC@{DC}!MatMgnt3d@{MatMgnt3d}}\index{MatMgnt3d@{MatMgnt3d}!DC@{DC}}\item[{\em 
DC\label{classMatMgnt3d_dca29a1140aadadfd92b34a02fa516ef_dca29a1140aadadfd92b34a02fa516efcf75e54791dd1f49f918345fdfe2430b}
}]\end{description}
\end{Desc}



Definition at line 41 of file matmgnt3d.hpp.

\footnotesize\begin{verbatim}41 {        UE=0,   UC=1,   DE=2,   DC=3,  };
\end{verbatim}\normalsize 




\subsection{Constructor \& Destructor Documentation}
\index{MatMgnt3d@{Mat\-Mgnt3d}!MatMgnt3d@{MatMgnt3d}}
\index{MatMgnt3d@{MatMgnt3d}!MatMgnt3d@{Mat\-Mgnt3d}}
\subsubsection{\setlength{\rightskip}{0pt plus 5cm}Mat\-Mgnt3d::Mat\-Mgnt3d ()}\label{classMatMgnt3d_f0b4e20a31097a5ff1315875944b0de9_f0b4e20a31097a5ff1315875944b0de9}




Definition at line 28 of file matmgnt3d.cpp.

\footnotesize\begin{verbatim}28                     : _np(6)
29 {
30 }
\end{verbatim}\normalsize 


\index{MatMgnt3d@{Mat\-Mgnt3d}!~MatMgnt3d@{$\sim$MatMgnt3d}}
\index{~MatMgnt3d@{$\sim$MatMgnt3d}!MatMgnt3d@{Mat\-Mgnt3d}}
\subsubsection{\setlength{\rightskip}{0pt plus 5cm}Mat\-Mgnt3d::$\sim$Mat\-Mgnt3d ()}\label{classMatMgnt3d_f6b631909bd602076d7b4dcbc5b9a8c4_f6b631909bd602076d7b4dcbc5b9a8c4}




Definition at line 32 of file matmgnt3d.cpp.

\footnotesize\begin{verbatim}33 {
34 }
\end{verbatim}\normalsize 




\subsection{Member Function Documentation}
\index{MatMgnt3d@{Mat\-Mgnt3d}!alt@{alt}}
\index{alt@{alt}!MatMgnt3d@{Mat\-Mgnt3d}}
\subsubsection{\setlength{\rightskip}{0pt plus 5cm}double Mat\-Mgnt3d::alt ()}\label{classMatMgnt3d_d3e453769e9e986f45d4fce5f3f4b9ad_d3e453769e9e986f45d4fce5f3f4b9ad}




Definition at line 36 of file matmgnt3d.cpp.

References \_\-np.

Referenced by Dwn\-Chk2Dwn\-Equ\_\-dgemv(), and Upw\-Chk2Upw\-Equ\_\-dgemv().

\footnotesize\begin{verbatim}37 {
38   return pow(0.1, _np+1);
39 }
\end{verbatim}\normalsize 


\index{MatMgnt3d@{Mat\-Mgnt3d}!cptwvv@{cptwvv}}
\index{cptwvv@{cptwvv}!MatMgnt3d@{Mat\-Mgnt3d}}
\subsubsection{\setlength{\rightskip}{0pt plus 5cm}int Mat\-Mgnt3d::cptwvv (int, double, fftw\_\-complex $\ast$, int, fftw\_\-complex $\ast$, int, fftw\_\-complex $\ast$, int)\hspace{0.3cm}{\tt  [inline]}}\label{classMatMgnt3d_a81bf8ce3d0c283d3bf63cae00ae99e6_a81bf8ce3d0c283d3bf63cae00ae99e6}




Definition at line 445 of file matmgnt3d.cpp.

References \_\-mmvec.

\footnotesize\begin{verbatim}446 {
447   for(int i=0; i<n; i++) {
448          (*z)[0] += alpha * ( (*x)[0] * (*y)[0] - (*x)[1] * (*y)[1]);
449          (*z)[1] += alpha * ( (*x)[0] * (*y)[1] + (*x)[1] * (*y)[0]);
450          x = x + incrementX;
451          y = y + incrementY;
452          z = z + incrementZ;
453   }  //iC( PetscLogFlops( 10*n ) );
454   return 0;
455 }
\end{verbatim}\normalsize 


\index{MatMgnt3d@{Mat\-Mgnt3d}!dim@{dim}}
\index{dim@{dim}!MatMgnt3d@{Mat\-Mgnt3d}}
\subsubsection{\setlength{\rightskip}{0pt plus 5cm}int Mat\-Mgnt3d::dim ()\hspace{0.3cm}{\tt  [inline]}}\label{classMatMgnt3d_e3f9953141ca7d7525b58c0ff1ef6b1f_e3f9953141ca7d7525b58c0ff1ef6b1f}


dimension = 3 

Definition at line 78 of file matmgnt3d.hpp.

Referenced by Dwn\-Chk2Dwn\-Equ\_\-dgemv(), Dwn\-Equ2Dwn\-Chk\_\-dgemv(), local\-Pos(), reg\-Pos\-Cal(), sam\-Pos\-Cal(), Upw\-Chk2Upw\-Equ\_\-dgemv(), Upw\-Equ2Dwn\-Chk\_\-dgemv(), and Upw\-Equ2Upw\-Chk\_\-dgemv().

\footnotesize\begin{verbatim}78 { return 3; }
\end{verbatim}\normalsize 


\index{MatMgnt3d@{Mat\-Mgnt3d}!DwnChk2DwnEqu_dgemv@{DwnChk2DwnEqu\_\-dgemv}}
\index{DwnChk2DwnEqu_dgemv@{DwnChk2DwnEqu\_\-dgemv}!MatMgnt3d@{Mat\-Mgnt3d}}
\subsubsection{\setlength{\rightskip}{0pt plus 5cm}int Mat\-Mgnt3d::Dwn\-Chk2Dwn\-Equ\_\-dgemv (int {\em level}, const Dbl\-Num\-Vec \&, Dbl\-Num\-Vec \&)}\label{classMatMgnt3d_b3f2a5306c0edee2c9df2722feb84c6c_b3f2a5306c0edee2c9df2722feb84c6c}




Definition at line 159 of file matmgnt3d.cpp.

References \_\-dwn\-Chk2Dwn\-Equ, \_\-hom, \_\-knl, alt(), DC, DE, dim(), Kernel3d::kernel(), pln\-Dat\-Sze(), and sam\-Pos().

\footnotesize\begin{verbatim}160 {
161   DblNumMat& _DC2DE = (_hom==true) ? _dwnChk2DwnEqu[0]: _dwnChk2DwnEqu[l];
162   double R = (_hom==true) ? 1 : 1.0/pow(2.0,l);
163   //---------compute matrix
164   if(_DC2DE.m()==0) {    //cerr<<"DwnChk2DwnEqu compute"<<endl;
165          DblNumMat dd2c(plnDatSze(DC), plnDatSze(DE));
166          DblNumMat chkPos(dim(),samPos(DC).n());                clear(chkPos);   iC( daxpy(R, samPos(DC), chkPos) ); //scale
167          DblNumMat denPos(dim(),samPos(DE).n());                clear(denPos);   iC( daxpy(R, samPos(DE), denPos) ); //scale
168          
169          iC( _knl.kernel(denPos, denPos, chkPos, dd2c) );//matrix
170          _DC2DE.resize(plnDatSze(DE), plnDatSze(DC)); //_memused[2] += plndnenum()*dof()*plndncnum()*dof()*sizeof(double);
171          iC( pinv(dd2c, alt(), _DC2DE) );
172   }
173   //---------matvec
174   if(_hom==true) {
175          int srcDOF = this->srcDOF();
176          DblNumVec tmpDen(srcDOF*samPos(DE).n(), false, _wsbuf);         clear(tmpDen);
177          iC( dgemv(1.0, _DC2DE, chk, 1.0, tmpDen) );
178          //scale
179          vector<double> sclvec(srcDOF);  for(int s=0; s<srcDOF; s++)            sclvec[s] = pow(2.0, - l*_degVec[s]);
180          int cnt = 0;
181          for(int i=0; i<samPos(DE).n(); i++)
182                 for(int s=0; s<srcDOF; s++) {
183                   den(cnt) = den(cnt) + tmpDen(cnt) * sclvec[s];
184                   cnt++;
185                 }
186   } else {
187          iC( dgemv(1.0, _DC2DE, chk, 1.0, den) );
188   }
189   return 0;
190 }
\end{verbatim}\normalsize 


\index{MatMgnt3d@{Mat\-Mgnt3d}!DwnEqu2DwnChk_dgemv@{DwnEqu2DwnChk\_\-dgemv}}
\index{DwnEqu2DwnChk_dgemv@{DwnEqu2DwnChk\_\-dgemv}!MatMgnt3d@{Mat\-Mgnt3d}}
\subsubsection{\setlength{\rightskip}{0pt plus 5cm}int Mat\-Mgnt3d::Dwn\-Equ2Dwn\-Chk\_\-dgemv (int {\em level}, Index3 {\em ii}, const Dbl\-Num\-Vec \&, Dbl\-Num\-Vec \&)}\label{classMatMgnt3d_87d8ed68b09cff332b257ad098ea81cc_87d8ed68b09cff332b257ad098ea81cc}




Definition at line 192 of file matmgnt3d.cpp.

References \_\-deg\-Vec, \_\-dwn\-Equ2Dwn\-Chk, \_\-hom, \_\-knl, \_\-wsbuf, DC, dim(), Kernel3d::kernel(), pln\-Dat\-Sze(), sam\-Pos(), and src\-DOF().

\footnotesize\begin{verbatim}193 {
194   NumTns<DblNumMat>& _DE2DC = (_hom==true) ? _dwnEqu2DwnChk[0] : _dwnEqu2DwnChk[l];
195   double R = (_hom==true) ? 1 : 1.0/pow(2.0, l);  //OffTns<DblNumMat>& _DwnEqu2DwnChk = _de2DwnChk[l];  double R       = 1.0/pow(2.0, l);
196   if(_DE2DC.m()==0)      _DE2DC.resize(2,2,2);
197   DblNumMat& _DE2DCii = _DE2DC(idx[0], idx[1], idx[2]);
198   
199   //---------compute matrix
200   if(_DE2DCii.m()==0) {  //cerr<<"DwnEqu2DwnChk compute"<<endl;
201          _DE2DCii.resize(plnDatSze(DC), plnDatSze(DE)); //_memused[3] += plndncnum()*dof()*plndnenum()*dof()*sizeof(double);
202          DblNumMat denPos(dim(),samPos(DE).n());                  clear(denPos);         iC( daxpy(R, samPos(DE), denPos) ); //scale
203          DblNumMat chkPos(dim(),samPos(DC).n());                  clear(chkPos);         iC( daxpy(0.5*R, samPos(DC), chkPos) ); //scale
204          for(int i=0; i<dim(); i++) for(int j=0; j<samPos(DC).n(); j++) chkPos(i,j) = chkPos(i,j) + (double(idx(i))-0.5)*R;
205          
206          iC( _knl.kernel(denPos, denPos, chkPos, _DE2DCii) );
207   }
208   //---------matvec
209   if(_hom==true) {
210          int srcDOF = this->srcDOF();
211          DblNumVec tmpDen(srcDOF*samPos(DE).n(), false, _wsbuf);         clear(tmpDen);
212          vector<double> sclvec(srcDOF);  for(int s=0; s<srcDOF; s++)            sclvec[s] = pow(2.0, l*_degVec[s]);
213          int cnt = 0;
214          for(int i=0; i<samPos(DE).n(); i++)
215                 for(int s=0; s<srcDOF; s++) {
216                   tmpDen(cnt) = den(cnt) * sclvec[s];
217                   cnt++;
218                 }
219          iC( dgemv(1.0, _DE2DCii, tmpDen, 1.0, chk) );
220   } else {
221          iC( dgemv(1.0, _DE2DCii, den, 1.0, chk) );
222   }
223   return (0);
224 }
\end{verbatim}\normalsize 


\index{MatMgnt3d@{Mat\-Mgnt3d}!effDatSze@{effDatSze}}
\index{effDatSze@{effDatSze}!MatMgnt3d@{Mat\-Mgnt3d}}
\subsubsection{\setlength{\rightskip}{0pt plus 5cm}int Mat\-Mgnt3d::eff\-Dat\-Sze (int {\em tp})}\label{classMatMgnt3d_2ce9326034b7f6cfc6196785759829e1_2ce9326034b7f6cfc6196785759829e1}


the size of eff data (data stored when using FFT) 

Definition at line 79 of file matmgnt3d.cpp.

References \_\-np, DE, src\-DOF(), trg\-DOF(), and UE.

\footnotesize\begin{verbatim}80 {
81   //int et = eq().et();
82   int effNum = (2*_np+2)*(2*_np)*(2*_np);
83   if(tp==UE || tp==DE)
84          return effNum*srcDOF();
85   else
86          return effNum*trgDOF();
87 }
\end{verbatim}\normalsize 


\index{MatMgnt3d@{Mat\-Mgnt3d}!effVal2PlnVal@{effVal2PlnVal}}
\index{effVal2PlnVal@{effVal2PlnVal}!MatMgnt3d@{Mat\-Mgnt3d}}
\subsubsection{\setlength{\rightskip}{0pt plus 5cm}int Mat\-Mgnt3d::eff\-Val2Pln\-Val (int {\em level}, const Dbl\-Num\-Vec \&, Dbl\-Num\-Vec \&)}\label{classMatMgnt3d_17b8e15860d967b5e69c356d0b11a215_17b8e15860d967b5e69c356d0b11a215}


effective-$>$regular-$>$plain 



Definition at line 277 of file matmgnt3d.cpp.

References \_\-np, reg\-Pos(), reg\-Val2Sam\-Val(), and trg\-DOF().

\footnotesize\begin{verbatim}278 {
279   DblNumVec regVal(regPos().n()*trgDOF());
280   int nnn[3];  nnn[0] = 2*_np;  nnn[1] = 2*_np;  nnn[2] = 2*_np;
281   fftw_plan invplan = fftw_plan_many_dft_c2r(3,nnn,trgDOF(), (fftw_complex*)(effVal.data()),NULL, trgDOF(),1, regVal.data(),NULL, trgDOF(),1, FFTW_ESTIMATE);
282   fftw_execute(invplan);
283   fftw_destroy_plan(invplan);  //rfftwnd_complex_to_real(_invplan, trgDOF(), (fftw_complex*)(effVal.data()), trgDOF(), 1, regVal.data(), trgDOF(), 1);
284   
285   iC( regVal2SamVal(regVal, plnVal) );
286   return (0);
287 }
\end{verbatim}\normalsize 


\index{MatMgnt3d@{Mat\-Mgnt3d}!getmmptr@{getmmptr}}
\index{getmmptr@{getmmptr}!MatMgnt3d@{Mat\-Mgnt3d}}
\subsubsection{\setlength{\rightskip}{0pt plus 5cm}{\bf Mat\-Mgnt3d} $\ast$ Mat\-Mgnt3d::getmmptr ({\bf Kernel3d}, int)\hspace{0.3cm}{\tt  [static]}}\label{classMatMgnt3d_79ba0648e2167456e6cd634d046ab0d6_79ba0648e2167456e6cd634d046ab0d6}




Definition at line 463 of file matmgnt3d.cpp.

References \_\-mmvec, knl(), np(), and setup().

\footnotesize\begin{verbatim}464 {
465   for(int i=0; i<_mmvec.size(); i++)
466          if(_mmvec[i].knl()==knl && _mmvec[i].np()==np)
467                 return &(_mmvec[i]);
468   
469   _mmvec.push_back( MatMgnt3d() );
470   int last = _mmvec.size()-1;
471   MatMgnt3d* tmp = &(_mmvec[last]); //get the last one
472   tmp->knl() = knl;  tmp->np() = np;
473   tmp->setup();
474   return tmp;
475 }
\end{verbatim}\normalsize 


\index{MatMgnt3d@{Mat\-Mgnt3d}!knl@{knl}}
\index{knl@{knl}!MatMgnt3d@{Mat\-Mgnt3d}}
\subsubsection{\setlength{\rightskip}{0pt plus 5cm}{\bf Kernel3d}\& Mat\-Mgnt3d::knl ()\hspace{0.3cm}{\tt  [inline]}}\label{classMatMgnt3d_d01911fc223ab4bc9604db12f942194e_d01911fc223ab4bc9604db12f942194e}




Definition at line 69 of file matmgnt3d.hpp.

References \_\-knl.

Referenced by getmmptr().

\footnotesize\begin{verbatim}69 { return _knl; }
\end{verbatim}\normalsize 


\index{MatMgnt3d@{Mat\-Mgnt3d}!localPos@{localPos}}
\index{localPos@{localPos}!MatMgnt3d@{Mat\-Mgnt3d}}
\subsubsection{\setlength{\rightskip}{0pt plus 5cm}int Mat\-Mgnt3d::local\-Pos (int, Point3, double, Dbl\-Num\-Mat \&)}\label{classMatMgnt3d_8b933eb3b49b51a202cfc077e84dca67_8b933eb3b49b51a202cfc077e84dca67}


return local position 

Definition at line 387 of file matmgnt3d.cpp.

References dim(), and sam\-Pos().

Referenced by FMM3d::Dwn\-Equ2Trg\-Chk\_\-dgemv(), FMM3d::Src\-Equ2Dwn\-Chk\_\-dgemv(), FMM3d::Src\-Equ2Upw\-Chk\_\-dgemv(), and FMM3d::Upw\-Equ2Trg\-Chk\_\-dgemv().

\footnotesize\begin{verbatim}388 {
389   const DblNumMat& bas = samPos(tp);
390   positions.resize(dim(), bas.n());
391   for(int i=0; i<dim(); i++)
392          for(int j=0; j<positions.n(); j++)
393                 positions(i,j) = center(i) + radius * bas(i,j);
394   return (0);
395 }
\end{verbatim}\normalsize 


\index{MatMgnt3d@{Mat\-Mgnt3d}!np@{np}}
\index{np@{np}!MatMgnt3d@{Mat\-Mgnt3d}}
\subsubsection{\setlength{\rightskip}{0pt plus 5cm}int\& Mat\-Mgnt3d::np ()\hspace{0.3cm}{\tt  [inline]}}\label{classMatMgnt3d_79545716c1fbc95cbd811a65ffa9f767_79545716c1fbc95cbd811a65ffa9f767}




Definition at line 70 of file matmgnt3d.hpp.

References \_\-np.

Referenced by getmmptr(), reg\-Val2Sam\-Val(), and sam\-Den2Reg\-Den().

\footnotesize\begin{verbatim}70 { return _np; }
\end{verbatim}\normalsize 


\index{MatMgnt3d@{Mat\-Mgnt3d}!plnDatSze@{plnDatSze}}
\index{plnDatSze@{plnDatSze}!MatMgnt3d@{Mat\-Mgnt3d}}
\subsubsection{\setlength{\rightskip}{0pt plus 5cm}int Mat\-Mgnt3d::pln\-Dat\-Sze (int {\em tp})}\label{classMatMgnt3d_7b6cc287face26f0d82bc4676e69fb83_7b6cc287face26f0d82bc4676e69fb83}


the size of plainn data 

Definition at line 71 of file matmgnt3d.cpp.

References \_\-sam\-Pos, DE, src\-DOF(), trg\-DOF(), and UE.

Referenced by FMM3d::dat\-Sze(), Dwn\-Chk2Dwn\-Equ\_\-dgemv(), Dwn\-Equ2Dwn\-Chk\_\-dgemv(), Upw\-Chk2Upw\-Equ\_\-dgemv(), and Upw\-Equ2Upw\-Chk\_\-dgemv().

\footnotesize\begin{verbatim}72 {
73   if(tp==UE || tp==DE)
74          return _samPos[tp].n()*srcDOF();
75   else
76          return _samPos[tp].n()*trgDOF();
77 }
\end{verbatim}\normalsize 


\index{MatMgnt3d@{Mat\-Mgnt3d}!plnDen2EffDen@{plnDen2EffDen}}
\index{plnDen2EffDen@{plnDen2EffDen}!MatMgnt3d@{Mat\-Mgnt3d}}
\subsubsection{\setlength{\rightskip}{0pt plus 5cm}int Mat\-Mgnt3d::pln\-Den2Eff\-Den (int {\em level}, const Dbl\-Num\-Vec \&, Dbl\-Num\-Vec \&)}\label{classMatMgnt3d_401736fd8ff6ae54a15badf8b0a91e33_401736fd8ff6ae54a15badf8b0a91e33}


plain-$>$regular-$>$effective 



Definition at line 226 of file matmgnt3d.cpp.

References \_\-deg\-Vec, \_\-hom, \_\-wsbuf, reg\-Pos(), sam\-Den2Reg\-Den(), sam\-Pos(), and src\-DOF().

\footnotesize\begin{verbatim}227 {
228   DblNumVec regDen(regPos().n()*srcDOF()); clear(regDen);
229   //iC( samDen2RegularDen(plnDen, regDen) );
230   //rfftwnd_real_to_complex(_forplan, srcDOF(), regDen.data(), srcDOF(), 1, (fftw_complex*)(effDen.data()), srcDOF(), 1);
231   if(_hom==true) {
232          int srcDOF = this->srcDOF();
233          DblNumVec tmpDen(srcDOF*samPos(UE).n(), false, _wsbuf);         clear(tmpDen);
234          vector<double> sclvec(srcDOF);  for(int s=0; s<srcDOF; s++)            sclvec[s] = pow(2.0, l*_degVec[s]);
235          int cnt = 0;
236          for(int i=0; i<samPos(UE).n(); i++)
237                 for(int s=0; s<srcDOF; s++) {
238                   tmpDen(cnt) = plnDen(cnt) * sclvec[s];
239                   cnt++;
240                 }
241          iC( samDen2RegDen(tmpDen, regDen) );
242   } else {
243          iC( samDen2RegDen(plnDen, regDen) );
244   }
245   
246   int nnn[3];  nnn[0] = 2*_np;  nnn[1] = 2*_np;  nnn[2] = 2*_np;
247   fftw_plan forplan = fftw_plan_many_dft_r2c(3,nnn,srcDOF(), regDen.data(),NULL, srcDOF(),1, (fftw_complex*)(effDen.data()),NULL, srcDOF(),1, FFTW_ESTIMATE);
248   fftw_execute(forplan);
249   fftw_destroy_plan(forplan);
250   //rfftwnd_real_to_complex(_forplan, srcDOF(), regDen.data(), srcDOF(), 1, (fftw_complex*)(effDen.data()), srcDOF(), 1);
251   
252   return (0);
253 }
\end{verbatim}\normalsize 


\index{MatMgnt3d@{Mat\-Mgnt3d}!regPos@{regPos}}
\index{regPos@{regPos}!MatMgnt3d@{Mat\-Mgnt3d}}
\subsubsection{\setlength{\rightskip}{0pt plus 5cm}const Dbl\-Num\-Mat\& Mat\-Mgnt3d::reg\-Pos ()\hspace{0.3cm}{\tt  [inline]}}\label{classMatMgnt3d_3e4e5fb8bf4f06cf8295d32216124585_3e4e5fb8bf4f06cf8295d32216124585}


return regular positions 

Definition at line 108 of file matmgnt3d.hpp.

References \_\-reg\-Pos.

Referenced by eff\-Val2Pln\-Val(), pln\-Den2Eff\-Den(), and Upw\-Equ2Dwn\-Chk\_\-dgemv().

\footnotesize\begin{verbatim}108 { return _regPos; }
\end{verbatim}\normalsize 


\index{MatMgnt3d@{Mat\-Mgnt3d}!regPosCal@{regPosCal}}
\index{regPosCal@{regPosCal}!MatMgnt3d@{Mat\-Mgnt3d}}
\subsubsection{\setlength{\rightskip}{0pt plus 5cm}int Mat\-Mgnt3d::reg\-Pos\-Cal (int {\em n}, double {\em R}, Dbl\-Num\-Mat \& {\em ret})}\label{classMatMgnt3d_4169f2bfee02545e9e3579c703129d1f_4169f2bfee02545e9e3579c703129d1f}


calculate the fft regular position grid 

Definition at line 423 of file matmgnt3d.cpp.

References dim().

Referenced by setup().

\footnotesize\begin{verbatim}424 {
425   int n = 2*np*2*np*2*np;
426   positions.resize(dim(), n);
427   double step = 2.0/(np-1);
428   int cnt = 0;
429   for(int k=0; k<2*np; k++)
430          for(int j=0; j<2*np; j++)
431                 for(int i=0; i<2*np; i++) {
432                   int gi = (i<np) ? i : i-2*np;
433                   int gj = (j<np) ? j : j-2*np;
434                   int gk = (k<np) ? k : k-2*np;
435                   positions(0, cnt) = R * gi*step;
436                   positions(1, cnt) = R * gj*step;
437                   positions(2, cnt) = R * gk*step;
438                   cnt ++;
439                 }
440   iA(cnt==n);
441   return 0;
442 }
\end{verbatim}\normalsize 


\index{MatMgnt3d@{Mat\-Mgnt3d}!regVal2SamVal@{regVal2SamVal}}
\index{regVal2SamVal@{regVal2SamVal}!MatMgnt3d@{Mat\-Mgnt3d}}
\subsubsection{\setlength{\rightskip}{0pt plus 5cm}int Mat\-Mgnt3d::reg\-Val2Sam\-Val (const Dbl\-Num\-Vec \&, Dbl\-Num\-Vec \&)}\label{classMatMgnt3d_157f482e203a0f8c33570e8f0fbccdc4_157f482e203a0f8c33570e8f0fbccdc4}




Definition at line 289 of file matmgnt3d.cpp.

References \_\-np, np(), and trg\-DOF().

Referenced by eff\-Val2Pln\-Val().

\footnotesize\begin{verbatim}290 {
291   int np = _np;
292   int rgnum = 2*np;
293   int trgDOF = this->trgDOF();
294   int cnt=0;
295   //the order of iterating is the same as SampleGrid
296   for(int i=0; i<np; i++)
297          for(int j=0; j<np; j++)
298                 for(int k=0; k<np; k++) {
299                   if(i==0 || i==np-1 || j==0 || j==np-1 || k==0 || k==np-1) {
300                          //the position is fortran style
301                          int rgoff = (k+np/2)*rgnum*rgnum + (j+np/2)*rgnum + (i+np/2);
302                          for(int f=0; f<trgDOF; f++) {
303                                 samVal(trgDOF*cnt + f) += regVal(trgDOF*rgoff + f);
304                          }
305                          cnt++;
306                   }
307                 }
308   return 0;
309 }
\end{verbatim}\normalsize 


\index{MatMgnt3d@{Mat\-Mgnt3d}!report@{report}}
\index{report@{report}!MatMgnt3d@{Mat\-Mgnt3d}}
\subsubsection{\setlength{\rightskip}{0pt plus 5cm}int Mat\-Mgnt3d::report ()}\label{classMatMgnt3d_ae2b28144bd345f29a29ecd1b43d62b7_ae2b28144bd345f29a29ecd1b43d62b7}




Definition at line 59 of file matmgnt3d.cpp.

References \_\-dwn\-Chk2Dwn\-Equ, \_\-dwn\-Equ2Dwn\-Chk, \_\-upw\-Chk2Upw\-Equ, \_\-upw\-Equ2Dwn\-Chk, and \_\-upw\-Equ2Upw\-Chk.

Referenced by FMM3d::evaluate().

\footnotesize\begin{verbatim}60 {
61   cerr << "matrix map size"<<endl;
62   cerr << _upwEqu2UpwChk.size() << " ";
63   cerr << _upwChk2UpwEqu.size() << " ";
64   cerr << _dwnChk2DwnEqu.size() << " ";
65   cerr << _dwnEqu2DwnChk.size() << " ";
66   cerr << _upwEqu2DwnChk.size() << endl;
67   return (0);
68 }
\end{verbatim}\normalsize 


\index{MatMgnt3d@{Mat\-Mgnt3d}!samDen2RegDen@{samDen2RegDen}}
\index{samDen2RegDen@{samDen2RegDen}!MatMgnt3d@{Mat\-Mgnt3d}}
\subsubsection{\setlength{\rightskip}{0pt plus 5cm}int Mat\-Mgnt3d::sam\-Den2Reg\-Den (const Dbl\-Num\-Vec \&, Dbl\-Num\-Vec \&)}\label{classMatMgnt3d_fa20d80fbb117bd78c6bad8de4964df7_fa20d80fbb117bd78c6bad8de4964df7}




Definition at line 255 of file matmgnt3d.cpp.

References \_\-np, np(), and src\-DOF().

Referenced by pln\-Den2Eff\-Den().

\footnotesize\begin{verbatim}256 {
257   int np = _np;
258   int rgnum = 2*np;
259   int srcDOF = this->srcDOF();
260   int cnt=0;
261   //the order of iterating is the same as SampleGrid
262   for(int i=0; i<np; i++)
263          for(int j=0; j<np; j++)
264                 for(int k=0; k<np; k++) {
265                   if(i==0 || i==np-1 || j==0 || j==np-1 || k==0 || k==np-1) {
266                          //the position is fortran style
267                          int rgoff = (k+np/2)*rgnum*rgnum + (j+np/2)*rgnum + (i+np/2);
268                          for(int f=0; f<srcDOF; f++) {
269                                 regDen(srcDOF*rgoff + f) += samDen(srcDOF*cnt + f);
270                          }
271                          cnt++;
272                   }
273                 }  //iC( PetscLogFlops(np*np*np*dof) );
274   return 0;
275 }
\end{verbatim}\normalsize 


\index{MatMgnt3d@{Mat\-Mgnt3d}!samPos@{samPos}}
\index{samPos@{samPos}!MatMgnt3d@{Mat\-Mgnt3d}}
\subsubsection{\setlength{\rightskip}{0pt plus 5cm}const Dbl\-Num\-Mat\& Mat\-Mgnt3d::sam\-Pos (int {\em tp})\hspace{0.3cm}{\tt  [inline]}}\label{classMatMgnt3d_86bc928d5146f28060644f8132bfbadd_86bc928d5146f28060644f8132bfbadd}


Return sample positions of UE=0,UC=1,DE=2, or DC=3 where UE = upward equivalent, UC = upward check, DE = downward equivalent, DC = downward check 

Definition at line 106 of file matmgnt3d.hpp.

References \_\-sam\-Pos.

Referenced by Dwn\-Chk2Dwn\-Equ\_\-dgemv(), Dwn\-Equ2Dwn\-Chk\_\-dgemv(), local\-Pos(), pln\-Den2Eff\-Den(), Upw\-Chk2Upw\-Equ\_\-dgemv(), and Upw\-Equ2Upw\-Chk\_\-dgemv().

\footnotesize\begin{verbatim}106 { return _samPos[tp]; }
\end{verbatim}\normalsize 


\index{MatMgnt3d@{Mat\-Mgnt3d}!samPosCal@{samPosCal}}
\index{samPosCal@{samPosCal}!MatMgnt3d@{Mat\-Mgnt3d}}
\subsubsection{\setlength{\rightskip}{0pt plus 5cm}int Mat\-Mgnt3d::sam\-Pos\-Cal (int {\em n}, double {\em R}, Dbl\-Num\-Mat \& {\em ret})}\label{classMatMgnt3d_632bf1ef6adf396e012a69cc9d71a202_632bf1ef6adf396e012a69cc9d71a202}


calculate various grids 

Definition at line 398 of file matmgnt3d.cpp.

References dim().

Referenced by setup().

\footnotesize\begin{verbatim}399 {
400   int n = np*np*np - (np-2)*(np-2)*(np-2);
401   positions.resize(dim(),n);
402   double step = 2.0/(np-1);
403   double init = -1.0;
404   int cnt = 0;
405   for(int i=0; i<np; i++)
406          for(int j=0; j<np; j++)
407                 for(int k=0; k<np; k++) {
408                   if(i==0 || i==np-1 || j==0 || j==np-1 || k==0 || k==np-1) {
409                          double x = init + i*step;
410                          double y = init + j*step;
411                          double z = init + k*step;
412                          positions(0,cnt) = R*x;
413                          positions(1,cnt) = R*y;
414                          positions(2,cnt) = R*z;
415                          cnt++;
416                   }
417                 }
418   iA(cnt==n);
419   return 0;
420 }
\end{verbatim}\normalsize 


\index{MatMgnt3d@{Mat\-Mgnt3d}!setup@{setup}}
\index{setup@{setup}!MatMgnt3d@{Mat\-Mgnt3d}}
\subsubsection{\setlength{\rightskip}{0pt plus 5cm}int Mat\-Mgnt3d::setup ()}\label{classMatMgnt3d_415266b005d267e801176f27c0434a50_415266b005d267e801176f27c0434a50}




Definition at line 42 of file matmgnt3d.cpp.

References \_\-deg\-Vec, \_\-hom, \_\-knl, \_\-np, \_\-reg\-Pos, \_\-sam\-Pos, DC, DE, Kernel3d::homogeneous(), Kernel3d::homogeneous\-Deg(), reg\-Pos\-Cal(), sam\-Pos\-Cal(), src\-DOF(), UC, and UE.

Referenced by getmmptr().

\footnotesize\begin{verbatim}43 {
44   //--------------------------------------------------------
45   _hom = _knl.homogeneous();
46   if(_hom==true) {
47          _knl.homogeneousDeg(_degVec); iA(_degVec.size()==srcDOF());
48   }
49   iC( samPosCal(_np,   1.0, _samPos[UE]) );
50   iC( samPosCal(_np+2, 3.0, _samPos[UC]) );
51   iC( samPosCal(_np,   3.0, _samPos[DE]) );
52   iC( samPosCal(_np,   1.0, _samPos[DC]) );
53   
54   iC( regPosCal(_np,   1.0, _regPos    ) ); //only one regPos
55   //--------------------------------------------------------
56   return (0);
57 }
\end{verbatim}\normalsize 


\index{MatMgnt3d@{Mat\-Mgnt3d}!srcDOF@{srcDOF}}
\index{srcDOF@{srcDOF}!MatMgnt3d@{Mat\-Mgnt3d}}
\subsubsection{\setlength{\rightskip}{0pt plus 5cm}int Mat\-Mgnt3d::src\-DOF ()\hspace{0.3cm}{\tt  [inline]}}\label{classMatMgnt3d_6b422ba557bb24eff82ac7125e206568_6b422ba557bb24eff82ac7125e206568}


src degree of freedom 

Definition at line 74 of file matmgnt3d.hpp.

References \_\-knl, and Kernel3d::src\-DOF().

Referenced by Dwn\-Equ2Dwn\-Chk\_\-dgemv(), eff\-Dat\-Sze(), pln\-Dat\-Sze(), pln\-Den2Eff\-Den(), sam\-Den2Reg\-Den(), setup(), and Upw\-Equ2Upw\-Chk\_\-dgemv().

\footnotesize\begin{verbatim}74 { return _knl.srcDOF(); }  //int tdof() { return eq().tdof(qt()); }
\end{verbatim}\normalsize 


\index{MatMgnt3d@{Mat\-Mgnt3d}!trgDOF@{trgDOF}}
\index{trgDOF@{trgDOF}!MatMgnt3d@{Mat\-Mgnt3d}}
\subsubsection{\setlength{\rightskip}{0pt plus 5cm}int Mat\-Mgnt3d::trg\-DOF ()\hspace{0.3cm}{\tt  [inline]}}\label{classMatMgnt3d_ea43c1f768349e86110e4e2d06e38d10_ea43c1f768349e86110e4e2d06e38d10}


trg degree of freedom 

Definition at line 76 of file matmgnt3d.hpp.

References \_\-knl, and Kernel3d::trg\-DOF().

Referenced by eff\-Dat\-Sze(), eff\-Val2Pln\-Val(), pln\-Dat\-Sze(), reg\-Val2Sam\-Val(), and Upw\-Equ2Dwn\-Chk\_\-dgemv().

\footnotesize\begin{verbatim}76 { return _knl.trgDOF(); }
\end{verbatim}\normalsize 


\index{MatMgnt3d@{Mat\-Mgnt3d}!UpwChk2UpwEqu_dgemv@{UpwChk2UpwEqu\_\-dgemv}}
\index{UpwChk2UpwEqu_dgemv@{UpwChk2UpwEqu\_\-dgemv}!MatMgnt3d@{Mat\-Mgnt3d}}
\subsubsection{\setlength{\rightskip}{0pt plus 5cm}int Mat\-Mgnt3d::Upw\-Chk2Upw\-Equ\_\-dgemv (int {\em level}, const Dbl\-Num\-Vec \&, Dbl\-Num\-Vec \&)}\label{classMatMgnt3d_6b9d3c2c722c190a882a82cef0429893_6b9d3c2c722c190a882a82cef0429893}




Definition at line 90 of file matmgnt3d.cpp.

References \_\-hom, \_\-knl, \_\-upw\-Chk2Upw\-Equ, alt(), dim(), Kernel3d::kernel(), pln\-Dat\-Sze(), sam\-Pos(), UC, and UE.

\footnotesize\begin{verbatim}91 {
92   DblNumMat& _UC2UE = (_hom==true) ? _upwChk2UpwEqu[0] : _upwChk2UpwEqu[l];
93   double R = (_hom==true) ? 1 : 1.0/pow(2.0,l);
94   //---------compute matrix
95   if(_UC2UE.m()==0) {    //cerr<<"UpwChk2UpwEqu compute"<<endl;
96          //set matrix
97          DblNumMat ud2c(plnDatSze(UC), plnDatSze(UE));
98          DblNumMat chkPos(dim(),samPos(UC).n());         clear(chkPos);  iC( daxpy(R, samPos(UC), chkPos) ); //scale
99          DblNumMat denPos(dim(),samPos(UE).n());         clear(denPos);  iC( daxpy(R, samPos(UE), denPos) ); //scale
100          
101          iC( _knl.kernel(denPos, denPos, chkPos, ud2c) );
102          _UC2UE.resize(plnDatSze(UE), plnDatSze(UC));
103          iC( pinv(ud2c, alt(), _UC2UE) );
104   }
105   //---------matvec
106   if(_hom==true) {
107          //matvec
108          int srcDOF = this->srcDOF();
109          DblNumVec tmpDen(srcDOF*samPos(UE).n(), false, _wsbuf);         clear(tmpDen);
110          iC( dgemv(1.0, _UC2UE, chk, 1.0, tmpDen) );
111          //scale
112          vector<double> sclvec(srcDOF);  for(int s=0; s<srcDOF; s++)            sclvec[s] = pow(2.0, - l*_degVec[s]);
113          int cnt = 0;
114          for(int i=0; i < samPos(UE).n(); i++)
115                 for(int s=0; s < srcDOF; s++) {
116                   den(cnt) = den(cnt) + tmpDen(cnt) * sclvec[s];
117                   cnt++;
118                 }
119   } else {
120          iC( dgemv(1.0, _UC2UE, chk, 1.0, den) );
121   }
122   return (0);
123 }
\end{verbatim}\normalsize 


\index{MatMgnt3d@{Mat\-Mgnt3d}!UpwEqu2DwnChk_dgemv@{UpwEqu2DwnChk\_\-dgemv}}
\index{UpwEqu2DwnChk_dgemv@{UpwEqu2DwnChk\_\-dgemv}!MatMgnt3d@{Mat\-Mgnt3d}}
\subsubsection{\setlength{\rightskip}{0pt plus 5cm}int Mat\-Mgnt3d::Upw\-Equ2Dwn\-Chk\_\-dgemv (int {\em level}, Index3 {\em ii}, const Dbl\-Num\-Vec \& {\em eff\-Den}, Dbl\-Num\-Vec \& {\em eff\-Val})}\label{classMatMgnt3d_265501b782aca2ac9f06197ab087e319_265501b782aca2ac9f06197ab087e319}




Definition at line 311 of file matmgnt3d.cpp.

References \_\-hom, \_\-knl, \_\-np, \_\-upw\-Equ2Dwn\-Chk, dim(), Kernel3d::kernel(), reg\-Pos(), and trg\-DOF().

\footnotesize\begin{verbatim}312 {
313   OffTns<DblNumMat>& _UpwEqu2DwnChk = (_hom==true) ? _upwEqu2DwnChk[0] : _upwEqu2DwnChk[l];
314   double R = (_hom==true) ? 1.0 : 1.0/pow(2.0, l); //OffTns< DblNumMat >& _UpwEqu2DwnChk = _upwEqu2DwnChk[l];  double R       = 1.0/pow(2.0, l);
315   if(_UpwEqu2DwnChk.m()==0)      _UpwEqu2DwnChk.resize(7,7,7,-3,-3,-3);
316   //DblNumMat* _UpwEqu2DwnChkii = new DblNumMat;
317   DblNumMat& _UpwEqu2DwnChkii = _UpwEqu2DwnChk(idx[0], idx[1], idx[2]);
318   
319   int effNum = (2*_np+2)*(2*_np)*(2*_np);
320   int srcDOF = this->srcDOF();
321   int trgDOF = this->trgDOF();
322   
323   //---------compute matrix
324   if(_UpwEqu2DwnChkii.m()==0) { //compute it if necessary
325          //-----------------------       //cerr<<"UpwEqu2DwnChk(FFT) compute"<<endl;     //COMPUTE FFT   //Index3 idx = ii.idx();        
326          iA( idx.linfty()>1 );
327          DblNumMat denPos(dim(),1);      for(int i=0; i<dim(); i++)             denPos(i,0) = double(idx(i))*2.0*R; //shift
328          DblNumMat chkPos(dim(),regPos().n());   clear(chkPos);  iC( daxpy(R, regPos(), chkPos) );
329          DblNumMat tt(regPos().n()*trgDOF, srcDOF);
330          iC( _knl.kernel(denPos, denPos, chkPos, tt) );
331          // move data to tmp
332          DblNumMat tmp(trgDOF,regPos().n()*srcDOF);
333          for(int k=0; k<regPos().n();k++) {
334                 for(int i=0; i<trgDOF; i++)
335                   for(int j=0; j<srcDOF; j++) {
336                          tmp(i,j+k*srcDOF) = tt(i+k*trgDOF,j);
337                   }
338          }
339          _UpwEqu2DwnChkii.resize(trgDOF*srcDOF, effNum); //_memused[4] += dof*dof*effNum(UE)*sizeof(double);
340          //forward FFT from tmp to _UpwEqu2DwnChkii;
341          
342          int nnn[3];  nnn[0] = 2*_np;  nnn[1] = 2*_np;  nnn[2] = 2*_np;
343          fftw_plan forplan = fftw_plan_many_dft_r2c(3,nnn,srcDOF*trgDOF, tmp.data(),NULL, srcDOF*trgDOF,1, (fftw_complex*)(_UpwEqu2DwnChkii.data()),NULL, srcDOF*trgDOF,1, FFTW_ESTIMATE);
344          fftw_execute(forplan);
345          fftw_destroy_plan(forplan);
346 
347          //rfftwnd_real_to_complex(_forplan, srcDOF*trgDOF, tmp.data(), srcDOF*trgDOF, 1, (fftw_complex*)(_UpwEqu2DwnChkii.data()), srcDOF*trgDOF, 1);
348   }
349   //---------matvec
350   /*
351   //scale
352   DblNumVec tmpDen(effDen.m(), false, _wsbuf);  clear(tmpDen);
353   fftw_complex* denptr = NULL;// = (fftw_complex*)(tmpDen.data());
354   if(_hom==true) {
355          vector<double> sclvec(srcDOF);  for(int s=0; s<srcDOF; s++)            sclvec[s] = pow(2.0, l*_degVec[s]);
356          int cnt = 0;
357          for(int i=0; i<effDen.m()/srcDOF; i++)
358                 for(int s=0; s<srcDOF; s++) {
359                   tmpDen(cnt) = effDen(cnt) * sclvec[s];                cnt++;
360                 }
361          denptr = (fftw_complex*)(tmpDen.data());
362   } else {
363          denptr = (fftw_complex*)(effDen.data());
364   }
365   */
366   //fft mult
367   double nrmfc = 1.0/double(regPos().n());
368   fftw_complex* matPtr  = (fftw_complex*)(_UpwEqu2DwnChkii.data());
369   fftw_complex* denPtr = (fftw_complex*)(effDen.data());
370   fftw_complex* chkPtr   = (fftw_complex*)(effVal.data());
371   int matStp  = srcDOF*trgDOF;
372   int denStp = srcDOF;
373   int chkStp   = trgDOF;
374   
375   double newalpha = nrmfc;
376   for(int i=0; i<trgDOF; i++)
377          for(int j=0; j<srcDOF; j++) {
378                 int matOff = j*trgDOF + i;
379                 int denOff = j;
380                 int chkOff = i;
381                 iC( cptwvv(effNum/2, newalpha, matPtr+matOff, matStp, denPtr+denOff, denStp, chkPtr+chkOff, chkStp) );
382          }
383   return (0);
384 }
\end{verbatim}\normalsize 


\index{MatMgnt3d@{Mat\-Mgnt3d}!UpwEqu2UpwChk_dgemv@{UpwEqu2UpwChk\_\-dgemv}}
\index{UpwEqu2UpwChk_dgemv@{UpwEqu2UpwChk\_\-dgemv}!MatMgnt3d@{Mat\-Mgnt3d}}
\subsubsection{\setlength{\rightskip}{0pt plus 5cm}int Mat\-Mgnt3d::Upw\-Equ2Upw\-Chk\_\-dgemv (int {\em level}, Index3 {\em ii}, const Dbl\-Num\-Vec \&, Dbl\-Num\-Vec \&)}\label{classMatMgnt3d_c90ec518abbf6f391ef382697c60da9c_c90ec518abbf6f391ef382697c60da9c}




Definition at line 125 of file matmgnt3d.cpp.

References \_\-deg\-Vec, \_\-hom, \_\-knl, \_\-upw\-Equ2Upw\-Chk, \_\-wsbuf, dim(), Kernel3d::kernel(), pln\-Dat\-Sze(), sam\-Pos(), src\-DOF(), and UC.

\footnotesize\begin{verbatim}126 {
127   NumTns<DblNumMat>& _UE2UC = (_hom==true) ? _upwEqu2UpwChk[0] : _upwEqu2UpwChk[l];
128   double R = (_hom==true) ? 1 : 1.0/pow(2.0, l);  
129   if(_UE2UC.m()==0)      _UE2UC.resize(2,2,2);
130   DblNumMat& _UE2UCii = _UE2UC(idx(0), idx(1), idx(2));
131   //---------compute matrix
132   if(_UE2UCii.m()==0) {  //cerr<<"UpwEqu2UpwChk compute"<<endl;
133          _UE2UCii.resize(plnDatSze(UC), plnDatSze(UE)); //_memused[1] += plnnum(UC)*dof()*plnnum(UE)*dof()*sizeof(double);
134          DblNumMat chkPos(dim(),samPos(UC).n());         clear(chkPos);  iC( daxpy(2.0*R, samPos(UC), chkPos) ); //scale
135          DblNumMat denPos(dim(),samPos(UE).n());         clear(denPos);  iC( daxpy(R, samPos(UE), denPos) ); //scale
136          for(int i=0; i<dim(); i++) for(int j=0; j<samPos(UE).n(); j++) denPos(i,j) = denPos(i,j) + (2*idx(i)-1)*R;//shift
137          
138          iC( _knl.kernel(denPos, denPos, chkPos, _UE2UCii) );
139   }
140   //---------matvec
141   if(_hom==true) {
142          int srcDOF = this->srcDOF();
143          DblNumVec tmpDen(srcDOF*samPos(UE).n(), false, _wsbuf);         clear(tmpDen);
144          vector<double> sclvec(srcDOF);  for(int s=0; s<srcDOF; s++)            sclvec[s] = pow(2.0, l*_degVec[s]);
145          int cnt = 0;
146          for(int i=0; i<samPos(UE).n(); i++)
147                 for(int s=0; s<srcDOF; s++) {
148                   tmpDen(cnt) = den(cnt) * sclvec[s];
149                   cnt++;
150                 }
151          iC( dgemv(1.0, _UE2UCii, tmpDen, 1.0, chk) );
152   } else {
153          iC( dgemv(1.0, _UE2UCii, den, 1.0, chk) );
154   }
155   return (0);
156 }
\end{verbatim}\normalsize 




\subsection{Member Data Documentation}
\index{MatMgnt3d@{Mat\-Mgnt3d}!_degVec@{\_\-degVec}}
\index{_degVec@{\_\-degVec}!MatMgnt3d@{Mat\-Mgnt3d}}
\subsubsection{\setlength{\rightskip}{0pt plus 5cm}vector$<$double$>$ {\bf Mat\-Mgnt3d::\_\-deg\-Vec}}\label{classMatMgnt3d_7a8804dd4c063a2060da6494c8815b6b_7a8804dd4c063a2060da6494c8815b6b}




Definition at line 47 of file matmgnt3d.hpp.

Referenced by Dwn\-Equ2Dwn\-Chk\_\-dgemv(), pln\-Den2Eff\-Den(), setup(), and Upw\-Equ2Upw\-Chk\_\-dgemv().\index{MatMgnt3d@{Mat\-Mgnt3d}!_dwnChk2DwnEqu@{\_\-dwnChk2DwnEqu}}
\index{_dwnChk2DwnEqu@{\_\-dwnChk2DwnEqu}!MatMgnt3d@{Mat\-Mgnt3d}}
\subsubsection{\setlength{\rightskip}{0pt plus 5cm}map$<$int, Dbl\-Num\-Mat$>$ {\bf Mat\-Mgnt3d::\_\-dwn\-Chk2Dwn\-Equ}}\label{classMatMgnt3d_b9f0f9702035e0fc0ef0790827f5d059_b9f0f9702035e0fc0ef0790827f5d059}


Downward Check To Downward Equivalent 

Definition at line 55 of file matmgnt3d.hpp.

Referenced by Dwn\-Chk2Dwn\-Equ\_\-dgemv(), and report().\index{MatMgnt3d@{Mat\-Mgnt3d}!_dwnEqu2DwnChk@{\_\-dwnEqu2DwnChk}}
\index{_dwnEqu2DwnChk@{\_\-dwnEqu2DwnChk}!MatMgnt3d@{Mat\-Mgnt3d}}
\subsubsection{\setlength{\rightskip}{0pt plus 5cm}map$<$int, Num\-Tns$<$Dbl\-Num\-Mat$>$ $>$ {\bf Mat\-Mgnt3d::\_\-dwn\-Equ2Dwn\-Chk}}\label{classMatMgnt3d_9dda4a0f7503a36d6e846d77ca1289d7_9dda4a0f7503a36d6e846d77ca1289d7}


Downward Equivalent To Downward Check 

Definition at line 57 of file matmgnt3d.hpp.

Referenced by Dwn\-Equ2Dwn\-Chk\_\-dgemv(), and report().\index{MatMgnt3d@{Mat\-Mgnt3d}!_hom@{\_\-hom}}
\index{_hom@{\_\-hom}!MatMgnt3d@{Mat\-Mgnt3d}}
\subsubsection{\setlength{\rightskip}{0pt plus 5cm}bool {\bf Mat\-Mgnt3d::\_\-hom}}\label{classMatMgnt3d_d7665fc205b24f10a29dc5fb47e4c9f6_d7665fc205b24f10a29dc5fb47e4c9f6}




Definition at line 46 of file matmgnt3d.hpp.

Referenced by cpp2mex(), Dwn\-Chk2Dwn\-Equ\_\-dgemv(), Dwn\-Equ2Dwn\-Chk\_\-dgemv(), mex2cpp(), pln\-Den2Eff\-Den(), setup(), Upw\-Chk2Upw\-Equ\_\-dgemv(), Upw\-Equ2Dwn\-Chk\_\-dgemv(), and Upw\-Equ2Upw\-Chk\_\-dgemv().\index{MatMgnt3d@{Mat\-Mgnt3d}!_knl@{\_\-knl}}
\index{_knl@{\_\-knl}!MatMgnt3d@{Mat\-Mgnt3d}}
\subsubsection{\setlength{\rightskip}{0pt plus 5cm}{\bf Kernel3d} {\bf Mat\-Mgnt3d::\_\-knl}}\label{classMatMgnt3d_cb56a4f532a89bc1a51ee17c2434a106_cb56a4f532a89bc1a51ee17c2434a106}




Definition at line 44 of file matmgnt3d.hpp.

Referenced by cpp2mex(), Dwn\-Chk2Dwn\-Equ\_\-dgemv(), Dwn\-Equ2Dwn\-Chk\_\-dgemv(), knl(), mex2cpp(), setup(), src\-DOF(), trg\-DOF(), Upw\-Chk2Upw\-Equ\_\-dgemv(), Upw\-Equ2Dwn\-Chk\_\-dgemv(), and Upw\-Equ2Upw\-Chk\_\-dgemv().\index{MatMgnt3d@{Mat\-Mgnt3d}!_mmvec@{\_\-mmvec}}
\index{_mmvec@{\_\-mmvec}!MatMgnt3d@{Mat\-Mgnt3d}}
\subsubsection{\setlength{\rightskip}{0pt plus 5cm}vector$<$ {\bf Mat\-Mgnt3d} $>$ {\bf Mat\-Mgnt3d::\_\-mmvec}\hspace{0.3cm}{\tt  [static]}}\label{classMatMgnt3d_cb639bda4a8ef940adb1d3495123f208_cb639bda4a8ef940adb1d3495123f208}




Definition at line 119 of file matmgnt3d.hpp.

Referenced by cptwvv(), getmmptr(), and mex\-Function().\index{MatMgnt3d@{Mat\-Mgnt3d}!_np@{\_\-np}}
\index{_np@{\_\-np}!MatMgnt3d@{Mat\-Mgnt3d}}
\subsubsection{\setlength{\rightskip}{0pt plus 5cm}int {\bf Mat\-Mgnt3d::\_\-np}}\label{classMatMgnt3d_3d05941198fa7da648bc68ddd5d0b5c1_3d05941198fa7da648bc68ddd5d0b5c1}




Definition at line 45 of file matmgnt3d.hpp.

Referenced by alt(), cpp2mex(), eff\-Dat\-Sze(), eff\-Val2Pln\-Val(), mex2cpp(), np(), reg\-Val2Sam\-Val(), sam\-Den2Reg\-Den(), setup(), and Upw\-Equ2Dwn\-Chk\_\-dgemv().\index{MatMgnt3d@{Mat\-Mgnt3d}!_regPos@{\_\-regPos}}
\index{_regPos@{\_\-regPos}!MatMgnt3d@{Mat\-Mgnt3d}}
\subsubsection{\setlength{\rightskip}{0pt plus 5cm}Dbl\-Num\-Mat {\bf Mat\-Mgnt3d::\_\-reg\-Pos}}\label{classMatMgnt3d_afe2677ddc7bd25cc4255c723397bdd7_afe2677ddc7bd25cc4255c723397bdd7}


regular positions 

Definition at line 63 of file matmgnt3d.hpp.

Referenced by reg\-Pos(), and setup().\index{MatMgnt3d@{Mat\-Mgnt3d}!_samPos@{\_\-samPos}}
\index{_samPos@{\_\-samPos}!MatMgnt3d@{Mat\-Mgnt3d}}
\subsubsection{\setlength{\rightskip}{0pt plus 5cm}Dbl\-Num\-Mat {\bf Mat\-Mgnt3d::\_\-sam\-Pos}[4]}\label{classMatMgnt3d_ba2fc588d5e34912fb575182f6d110c4_ba2fc588d5e34912fb575182f6d110c4}


sample positions - different depending on whether UW, UC, DE or DC 

Definition at line 61 of file matmgnt3d.hpp.

Referenced by pln\-Dat\-Sze(), sam\-Pos(), and setup().\index{MatMgnt3d@{Mat\-Mgnt3d}!_upwChk2UpwEqu@{\_\-upwChk2UpwEqu}}
\index{_upwChk2UpwEqu@{\_\-upwChk2UpwEqu}!MatMgnt3d@{Mat\-Mgnt3d}}
\subsubsection{\setlength{\rightskip}{0pt plus 5cm}map$<$int, Dbl\-Num\-Mat$>$ {\bf Mat\-Mgnt3d::\_\-upw\-Chk2Upw\-Equ}}\label{classMatMgnt3d_b698e133e8b80af377013cfa83a64f9c_b698e133e8b80af377013cfa83a64f9c}


Upward Check To Upward Equivalent 

Definition at line 51 of file matmgnt3d.hpp.

Referenced by report(), and Upw\-Chk2Upw\-Equ\_\-dgemv().\index{MatMgnt3d@{Mat\-Mgnt3d}!_upwEqu2DwnChk@{\_\-upwEqu2DwnChk}}
\index{_upwEqu2DwnChk@{\_\-upwEqu2DwnChk}!MatMgnt3d@{Mat\-Mgnt3d}}
\subsubsection{\setlength{\rightskip}{0pt plus 5cm}map$<$int, Off\-Tns$<$Dbl\-Num\-Mat$>$ $>$ {\bf Mat\-Mgnt3d::\_\-upw\-Equ2Dwn\-Chk}}\label{classMatMgnt3d_4433042b9af17fab86b8f035b255d886_4433042b9af17fab86b8f035b255d886}


Upward Equivalent To Downward Check 

Definition at line 59 of file matmgnt3d.hpp.

Referenced by report(), and Upw\-Equ2Dwn\-Chk\_\-dgemv().\index{MatMgnt3d@{Mat\-Mgnt3d}!_upwEqu2UpwChk@{\_\-upwEqu2UpwChk}}
\index{_upwEqu2UpwChk@{\_\-upwEqu2UpwChk}!MatMgnt3d@{Mat\-Mgnt3d}}
\subsubsection{\setlength{\rightskip}{0pt plus 5cm}map$<$int, Num\-Tns$<$Dbl\-Num\-Mat$>$ $>$ {\bf Mat\-Mgnt3d::\_\-upw\-Equ2Upw\-Chk}}\label{classMatMgnt3d_99d61bd38f82cf6e8a44bdfe501e2dda_99d61bd38f82cf6e8a44bdfe501e2dda}


Upward Equivalent To Upward Check 

Definition at line 53 of file matmgnt3d.hpp.

Referenced by report(), and Upw\-Equ2Upw\-Chk\_\-dgemv().\index{MatMgnt3d@{Mat\-Mgnt3d}!_wsbuf@{\_\-wsbuf}}
\index{_wsbuf@{\_\-wsbuf}!MatMgnt3d@{Mat\-Mgnt3d}}
\subsubsection{\setlength{\rightskip}{0pt plus 5cm}double {\bf Mat\-Mgnt3d::\_\-wsbuf}\hspace{0.3cm}{\tt  [static]}}\label{classMatMgnt3d_526f8f4a77395e52ecb0eb6d09440edc_526f8f4a77395e52ecb0eb6d09440edc}




Definition at line 118 of file matmgnt3d.hpp.

Referenced by Dwn\-Equ2Dwn\-Chk\_\-dgemv(), pln\-Den2Eff\-Den(), and Upw\-Equ2Upw\-Chk\_\-dgemv().

The documentation for this class was generated from the following files:\begin{CompactItemize}
\item 
fmm3d/{\bf matmgnt3d.hpp}\item 
fmm3d/{\bf matmgnt3d.cpp}\end{CompactItemize}
